\clearpage\thispagestyle{empty}
\pdfbookmark[0]{Sources for Pallas (letters of 1 February and 24 April 1815)}{pallascorr1815-sources}
\begin{center}\textbf{Sources for Pallas (letters of 1 February and 24 April 1815)}\end{center}
\noindent This AI translation follows the two letters from Carl Friedrich Gauss to Johann Elert Bode under Pallas in \href{https://gdz.sub.uni-goettingen.de/id/PPN236007467}{\textit{Werke}, vol.~VI} (Göttingen, 1874), pp.~383--385. Bode printed the excerpts with a collective heading covering Pallas, the winter solstice of 1814, the comet of 1815 and barometric height measurement in the \textit{Astronomisches Jahrbuch für 1818} (Berlin, 1815), pp.~\href{https://api.digitale-sammlungen.de/iiif/image/v2/bsb10538347_00175/full/1600,/0/default.jpg}{167}--\href{https://api.digitale-sammlungen.de/iiif/image/v2/bsb10538347_00181/full/1600,/0/default.jpg}{173}. The links lead to the Bavarian State Library copy.
\par\smallskip\noindent The footnote attached to the mention of Enke’s ephemeris on p.~\href{https://api.digitale-sammlungen.de/iiif/image/v2/bsb10538347_00176/full/1600,/0/default.jpg}{168} reads, in translation, “It already appears above.” The section on barometric height measurement, with three tables and a worked example, runs from p.~\href{https://api.digitale-sammlungen.de/iiif/image/v2/bsb10538347_00177/full/1600,/0/default.jpg}{169} through the upper part of p.~\href{https://api.digitale-sammlungen.de/iiif/image/v2/bsb10538347_00180/full/1600,/0/default.jpg}{172} in the yearbook. The collected text in Volume~VI omits both.
\par\smallskip\noindent The yearbook gives 9~46~58 for the mean time on 25~March and 59~13~3 for the right ascension on 30~March. The collected letters give 9~46~38 and 59~14~3, respectively. The yearbook gives the longitude of perihelion as $146^\circ\,7'\,2''$, and the collected letters give $146^\circ\,7\,2$. The separate collected report of 8~April gives $246^\circ\,7'\,2''$. Translation and typesetting produced by AI.
\clearpage
